

\section{Le mythe de la "Version" comme lecture : Traduire n'est pas lire}

L'enseignement du grec ancien en France traverse une crise systémique qui dépasse largement la
simple question des effectifs ou des dotations horaires. Elle touche à l'essence même de la
transmission et, plus profondément, à la conception cognitive de l'acte de lire. Depuis des
décennies, la discipline repose sur un paradigme pédagogique ininterrogé : la centralité de la «
version » (traduction du grec vers le français) comme unique mode d'accès, de validation et
d'appropriation de la langue.\cite{II-3-1-ref1} Ce modèle, hérité d'une sédimentation historique
complexe mêlant tradition rhétorique et impératifs de sélection républicaine, postule implicitement
une équivalence fonctionnelle entre \textit{traduire} et \textit{lire}. Or, cette équivalence est un
mythe.

L'hypothèse centrale que nous défendons est que la «
version », telle qu'elle est pratiquée institutionnellement, relève d'une activité de résolution de
problèmes (\textit{Problem Solving}) qui entre en conflit direct avec les mécanismes neurocognitifs
de la lecture fluide (\textit{Information Processing}).\cite{II-3-1-ref3}

En nous appuyant sur les théories de l'automaticité 5, de la charge cognitive 6 et du modèle
neurobiologique Déclaratif/Procédural 7, nous démontrerons que la confusion entre ces deux régimes
cognitifs constitue le principal obstacle à l'acquisition d'une compétence réelle en grec ancien.
Cette analyse s'inscrit dans le sillage du rapport Charvet-Bauduin (2018), qui pointait déjà
l'urgence de refonder la didactique des langues anciennes pour enrayer l'effondrement des effectifs
et redonner du sens à cet enseignement.\cite{II-3-1-ref8}

Pour comprendre la persistance du modèle de la version, il est indispensable de retracer sa
généalogie. Loin d'être une méthode naturelle ou universelle, la prééminence de la traduction écrite
est le produit d'une histoire spécifique du système éducatif français, cristallisée à la fin du XIXe
siècle.



\subsection{De l'imprégnation rhétorique à la « gymnastique mentale »}

L'histoire de l'enseignement des langues anciennes en Occident a longtemps oscillé entre deux pôles
: une tradition d'imprégnation orale et de mémorisation, visant à former des orateurs capables de
penser et de s'exprimer dans la langue cible, et une tradition analytique, centrée sur la grammaire
et la logique.\cite{II-3-1-ref10} Jusqu'au XVIIIe siècle, le latin et, dans une moindre mesure, le
grec, étaient encore perçus comme des langues de communication savante. Cependant, le XIXe siècle
marque une rupture décisive avec l'avènement de la philologie scientifique et la réorganisation de
l'Instruction publique.

Sous l'impulsion de réformateurs comme Michel Bréal et Jules Ferry, l'enseignement secondaire
français des années 1880 opère une mutation profonde.\cite{II-3-1-ref11} Il ne s'agit plus de former
des locuteurs en langues anciennes, mais d'utiliser ces langues comme des instruments de formation
de l'esprit. Michel Bréal, dans ses conférences à la Sorbonne en 1891, théorise cette approche en
parlant de « gymnastique mentale ».\cite{II-3-1-ref13} La difficulté d'accès au texte grec, la
résistance qu'il oppose au déchiffrement immédiat, devient en soi une vertu pédagogique. L'effort
d'analyse grammaticale et de transposition lexicale est censé développer la rigueur logique, la
capacité d'abstraction et la maîtrise de la langue maternelle.\cite{II-3-1-ref10}

Dans cette perspective, la « version » n'est pas une fin en soi (comprendre ce que dit Thucydide),
mais un moyen de discipliner l'intelligence. Comme le notent les instructions officielles de
l'époque, et comme cela perdure dans l'inconscient scolaire, on apprend le grec pour mieux savoir le
français.\cite{II-3-1-ref16} Cette instrumentalisation a eu pour effet pervers de détourner
l'attention de la langue cible (le grec) vers la langue source (le français), installant durablement
l'idée que la compétence ultime de l'helléniste réside dans l'élégance de sa traduction française et
non dans sa fluidité de lecture en grec.



\subsection{La sanctuarisation de l'exercice de version}

Cette philosophie éducative s'est institutionnalisée à travers les grands concours de recrutement de
l'élite professorale : l'Agrégation et le Concours Général. Depuis leur rétablissement ou leur
réforme au XIXe siècle, ces concours ont fait de la version (et du thème) les épreuves
reines.\cite{II-3-1-ref17} La « version sèche » — traduction d'un texte court, complexe, hors
contexte, avec pour seule aide un dictionnaire bilingue — est devenue l'étalon-or de l'excellence
académique.\cite{II-3-1-ref2}

Ce modèle a ruisselé sur l'ensemble du système secondaire. Le « dogme des textes authentiques »,
critiqué par le rapport Charvet, impose aux élèves, dès les premières années d'apprentissage, la
confrontation avec des auteurs canoniques (Platon, Sophocle, Xénophon) dont la difficulté
linguistique est sans commune mesure avec leurs compétences réelles.\cite{II-3-1-ref1} Au lieu de
lire des textes adaptés à leur niveau pour développer progressivement leur compétence, les élèves
sont contraints de déchiffrer laborieusement des monuments littéraires, transformant chaque phrase
en une énigme grammaticale à résoudre.



\subsection{Le constat de faillite : le rapport charvet (2018)}

Le rapport remis en 2018 par Pascal Charvet et David Bauduin au ministre de l'Éducation nationale
dresse un bilan sans appel de cette tradition. Les auteurs constatent une érosion dramatique des
effectifs — une perte de plus de 40 000 élèves en neuf ans — et une crise des vocations
enseignantes.\cite{II-3-1-ref8} Plus grave encore, le rapport souligne l'inefficacité pédagogique du
modèle traditionnel pour le plus grand nombre : les élèves, après plusieurs années d'étude, restent
incapables de lire un texte simple sans béquilles.\cite{II-3-1-ref9}

Le rapport Charvet appelle à une « valorisation » qui passe par un changement de paradigme : sortir
de l'obsession de la traduction grammaticale pour renouer avec l'esprit des Humanités, la culture,
et une approche plus vivante de la langue. Il remet en cause le fétichisme du « texte authentique »
immédiat, rappelant que les Anciens eux-mêmes n'apprenaient pas leur langue en commençant par
l'Iliade, mais par des exercices progressifs.\cite{II-3-1-ref9} Ce document officiel marque une
reconnaissance institutionnelle, bien que tardive, que la version est devenue un obstacle
épistémologique à la lecture.

Pour démontrer en quoi la version s'oppose à la lecture, il convient d'abord de définir ce qu'est
l'acte de lire du point de vue de la psychologie cognitive et de la psycholinguistique. Loin d'être
un déchiffrement conscient, la lecture experte est un traitement de l'information
(\textit{Information Processing}) caractérisé par l'automaticité.



\subsection{La théorie de l'automaticité (laberge \& samuels)}

Le modèle théorique fondateur de la lecture fluide est la théorie de l'automaticité formulée par
LaBerge et Samuels (1974).\cite{II-3-1-ref5} Ce modèle repose sur un postulat économique du cerveau
: nos ressources attentionnelles sont limitées.

1. \textbf{Ressources Limitées :} Le lecteur dispose d'un « réservoir » fini d'attention consciente.

2. \textbf{Compétition Cognitive :} L'acte de lire implique deux niveaux de traitement : le bas
niveau (décodage des lettres, reconnaissance des mots) et le haut niveau (compréhension sémantique,
inférences, appréciation esthétique).

3. \textbf{Le Seuil d'Automaticité :} Pour que la compréhension (haut niveau) puisse advenir, le
décodage (bas niveau) doit être \textit{automatique}, c'est-à-dire s'effectuer sans consommation
d'attention consciente.\cite{II-3-1-ref20}

Si le lecteur doit consacrer son attention à identifier la morphologie d'un verbe ou à chercher le
sens d'un mot (comme c'est le cas dans la version), il sature sa mémoire de travail. Il ne reste
plus de ressources disponibles pour construire le sens global de la phrase ou du texte. C'est ce qui
explique pourquoi un élève peut avoir traduit correctement tous les mots d'une phrase grecque sans
en avoir compris le sens global : son énergie cognitive a été entièrement dissipée dans le micro-
traitement lexical et grammatical.\cite{II-3-1-ref22}



\subsection{La fluidité de lecture (\textit{reading fluency}) et le seuil lexical}

William Grabe et Paul Nation, figures majeures de la recherche en acquisition des langues secondes
(SLA), ont démontré que la fluidité (\textit{fluency}) n'est pas un luxe, mais une condition
\textit{sine qua non} de la compréhension.\cite{II-3-1-ref23}

Un lecteur fluide reconnaît les mots instantanément (« \textit{sight word recognition} ») et traite
le texte par larges unités de sens (chunks), et non mot à mot. Pour atteindre cet état, la recherche
a établi un seuil lexical précis : le lecteur doit connaître \textbf{95\% à 98\%} des mots du texte
qu'il lit.\cite{II-3-1-ref26}

\begin{itemize}\item \textbf{À 98\% de couverture :} La lecture est fluide, le plaisir est possible, et les 2\% de mots inconnus peuvent être devinés par le contexte (apprentissage incident).\item \textbf{À 80\% de couverture :} (ce qui correspond souvent à la situation d'un élève devant un texte authentique, même en maîtrisant les 1100 mots les plus fréquents 26), la compréhension est impossible sans aide extérieure. Le texte devient un mur infranchissable, forçant le basculement vers une stratégie de déchiffrement analytique.\end{itemize}

\subsection{Données oculométriques : la signature visuelle de la lecture}

Les études d'oculométrie (\textit{eye-tracking}) fournissent des preuves physiologiques de la
distinction entre lire et déchiffrer.

\begin{itemize}\item \textbf{Lecture Fluide :} Le regard progresse par saccades rapides et régulières vers la droite (pour les langues LTR), avec des fixations courtes (200-250 ms) et peu de régressions (retours en arrière).\cite{II-3-1-ref28} Le cerveau anticipe l'information (\textit{top-down processing}).\item \textbf{Lecture Laborieuse (ou Traduction) :} Les fixations sont longues, les saccades sont courtes, et les régressions sont massives.\cite{II-3-1-ref30} Le lecteur fait des allers-retours constants pour vérifier la morphologie ou reconsidérer un sens. Ce pattern visuel est la signature d'un traitement de l'information défaillant, où le cerveau lutte pour assembler les données.\end{itemize}Ainsi, la lecture véritable se définit par la vitesse, l'automaticité et la focalisation sur le
sens. C'est un processus d'intégration continue.

À l'opposé du flux continu de la lecture, la « version » scolaire s'apparente à une activité
cognitive radicalement différente : la résolution de problèmes (\textit{Problem Solving}). Cette
distinction, théorisée notamment par Wolfgang Lörscher, est essentielle pour comprendre
l'incompatibilité des deux démarches.



\subsection{La traduction comme rupture stratégique}

Lörscher (1991) définit la stratégie de traduction comme « une procédure potentiellement consciente
pour la solution d'un problème auquel un individu est confronté lorsqu'il traduit un segment de
texte ».\cite{II-3-1-ref31} Contrairement à la lecture qui vise l'inconscience des processus
(automaticité), la version exige la conscience des processus.

Dans l'exercice de version, l'élève opère en mode « \textit{Stop-and-Go} » :

1. Lecture d'un segment.

2. Arrêt face à une difficulté (mot inconnu, forme ambiguë).

3. Activation d'une stratégie de résolution (analyse des désinences, recherche dans le
dictionnaire).

4. Formulation d'une hypothèse de traduction.

5. Vérification et reprise.\cite{II-3-1-ref4}

Ce processus haché empêche la construction d'une représentation mentale cohérente du texte (le «
modèle de situation »). L'élève ne navigue pas dans le monde du texte (la Grèce antique, l'action
décrite) ; il navigue dans le monde des règles grammaticales et des équivalences lexicales. Il ne
lit pas l'histoire d'Achille ; il résout l'équation syntaxique de la phrase décrivant Achille.



\subsection{La théorie de la charge cognitive et l'effet d'attention divisée}

La Théorie de la Charge Cognitive (\textit{Cognitive Load Theory \- CLT}) de John Sweller offre un
cadre puissant pour analyser l'échec de la version comme méthode d'apprentissage.\cite{II-3-1-ref6}
La CLT distingue trois types de charge :

\begin{itemize}\item \textbf{Charge Intrinsèque :} La complexité inhérente au grec ancien (alphabet, flexions).\item \textbf{Charge Extrinsèque :} La charge imposée par la méthode d'enseignement (mauvaise conception du matériel, usage maladroit des outils).\item \textbf{Charge Germane :} La charge utile consacrée à l'apprentissage (création de schémas mentaux).\end{itemize}L'usage systématique du dictionnaire bilingue (le fameux Bailly ou Gaffiot) en cours de version
génère un Effet d'Attention Divisée (Split-Attention Effect) massif.\cite{II-3-1-ref35}

L'élève doit détourner son attention visuelle et cognitive du texte source pour aller chercher une
information dans le dictionnaire, puis revenir au texte. Cette rupture provoque une surcharge de la
mémoire de travail : l'information stockée (le début de la phrase) s'efface avant que l'information
nouvelle (le sens du mot trouvé) ne puisse y être intégrée.

Les études montrent que ce va-et-vient constant non seulement ralentit le processus, mais dégrade la
compréhension globale et l'acquisition du vocabulaire.\cite{II-3-1-ref30} L'élève apprend à
manipuler le dictionnaire, pas à lire la langue.



\subsection{Le piège de la traduction mentale}

Même sans dictionnaire, la version encourage la « traduction mentale » (mental translation), c'est-
à-dire le réflexe de traduire silencieusement en français pour comprendre le
grec.\cite{II-3-1-ref38}

Bien que la traduction mentale soit une étape naturelle chez les débutants, elle devient un frein si
elle perdure.\cite{II-3-1-ref28} Elle maintient le cerveau dans une dépendance à la L1 (langue
maternelle), empêchant l'accès direct au sens (direct access to meaning) dans la L2.

\begin{itemize}\item \textbf{Interférence Cognitive :} La traduction mentale ajoute une étape de traitement supplémentaire (recodage L2 \-\> L1), consommant des ressources précieuses et ralentissant la fluidité en dessous du seuil nécessaire à la compréhension globale.\cite{II-3-1-ref35}\item \textbf{Oculométrie :} Les données montrent que lors d'une tâche de traduction, les mouvements oculaires sont erratiques, témoignant d'une charge cognitive élevée et d'une difficulté à intégrer l'information.\cite{II-3-1-ref3}\end{itemize}\begin{table}[h!]\small\centering\begin{tabular}{| p{2.3cm} | p{3.5cm} | p{3.5cm} |}\hline\textbf{Caractéristique} & \textbf{Lecture (Information Processing)} & \textbf{Version (Problem Solving)} \\ \hline\textbf{Objectif} & Compréhension du sens (Message) & Production d'un texte cible (Code) \\ \hline\textbf{Processus} & Automatique, Fluide, Inconscient & Stratégique, Haché, Conscient \\ \hline\textbf{Mémoire de Travail} & Optimisée pour le sens global & Saturée par le décodage local \\ \hline\textbf{Mouvement Oculaire} & Saccades régulières, Anticipation & Fixations longues, Régressions \\ \hline\textbf{Rapport à la L1} & Accès direct au concept & Médiation systématique par la L1 \\ \hline\end{tabular}\end{table}Tableau 1 : Comparaison cognitive entre la lecture et la version.\cite{II-3-1-ref3}

L'incompatibilité entre version et lecture trouve son explication ultime dans l'architecture
neurobiologique du cerveau humain. Les travaux de Michael Ullman sur le modèle Déclaratif/Procédural
(DP Model) éclairent pourquoi la méthode « Grammaire-Traduction » échoue à produire des lecteurs
fluides.



\subsection{Le modèle dp (declarative/procedural) d'ullman}

Le modèle DP postule que l'apprentissage et l'usage du langage reposent sur deux systèmes de mémoire
cérébraux distincts, anatomiquement et fonctionnellement dissociés.\cite{II-3-1-ref7}



\subsection{1\. la mémoire déclarative (système temporal)}

\begin{itemize}\item \textbf{Siège :} Lobe temporal médian (incluant l'hippocampe).\item \textbf{Fonction :} Stockage des connaissances explicites (« savoir que »), des faits, des événements (mémoire épisodique) et du vocabulaire (lexique mental).\item \textbf{Caractéristiques :} Apprentissage rapide (on peut retenir un mot en une seule exposition), conscient, flexible. Mais la récupération de l'information est relativement lente et demande de l'attention.\cite{II-3-1-ref42}\item \textbf{Dans la Version :} C'est le système sur-sollicité par l'enseignement traditionnel. L'élève apprend des listes de vocabulaire et des règles de grammaire explicites (« \textit{le datif exprime le moyen} »). Il stocke des connaissances \textit{sur} la langue.\end{itemize}

\subsection{2\. la mémoire procédurale (système fronto-striatal)}

\begin{itemize}\item \textbf{Siège :} Ganglions de la base (striatum) et régions frontales (aire de Broca).\item \textbf{Fonction :} Apprentissage des habiletés motrices et cognitives (« savoir comment »), des séquences, et de la grammaire implicite (combinatoire syntaxique, morphologie régulière).\item \textbf{Caractéristiques :} Apprentissage lent (nécessite de nombreuses répétitions), inconscient, rigide. Mais une fois acquis, le traitement est \textbf{automatique}, ultra-rapide et sans effort.\cite{II-3-1-ref43}\item \textbf{Dans la Lecture :} C'est le système qui permet la fluidité. Le lecteur expert « sent » que la phrase est correcte ou comprend la structure syntaxique sans avoir besoin de nommer la règle grammaticale.\end{itemize}

\subsection{Le drame de la non-procéduralisation}

Le problème fondamental de la méthode de la version est qu'elle entraîne exclusivement la mémoire
déclarative, tout en espérant des résultats qui dépendent de la mémoire procédurale.

L'élève devient un expert en connaissances déclaratives (il sait réciter ses déclinaisons et
expliquer les règles de l'accentuation), mais ces connaissances ne se transfèrent pas
automatiquement vers le système procédural nécessaire à la lecture fluide.\cite{II-3-1-ref45}

Selon Ullman, bien que les adultes apprenant une L2 commencent souvent par le système déclaratif,
l'objectif de l'exposition à la langue est de permettre le basculement progressif vers le système
procédural.\cite{II-3-1-ref47} Or, la version, en maintenant l'élève dans une posture d'analyse
consciente permanente (monitoring), bloque ce transfert. Elle empêche l'automatisation en forçant le
cerveau à repasser systématiquement par la règle explicite.\cite{II-3-1-ref48}



\subsection{L'hypothèse du « monitor » de krashen revisitée}

Ce constat neurobiologique valide a posteriori l'Hypothèse du Moniteur de Stephen Krashen, formulée
dans les années 1980.\cite{II-3-1-ref48}

\begin{itemize}\item \textbf{Acquisition (Acquisition) :} Processus subconscient similaire à l'apprentissage de la langue maternelle, dépendant de l'input compréhensible. C'est ce qui construit la compétence fluide (système procédural).\item \textbf{Apprentissage (Learning) :} Connaissance consciente des règles (système déclaratif).\item \textbf{Le Rôle du Moniteur :} Le savoir conscient ne sert que de « moniteur » (vérificateur) pour corriger la production. Il ne génère pas la parole ou la compréhension fluide.\end{itemize}La version est, par essence, un exercice d'hyper-monitoring. Elle force l'élève à utiliser son
système déclaratif pour construire le sens, ce qui est cognitivement épuisant et inefficace pour la
lecture cursive. Elle crée des « savants de la grammaire » qui sont des « analphabètes fonctionnels
» devant un texte.\cite{II-3-1-ref51}

Face à l'impasse cognitive de la version, des alternatives pédagogiques émergent, fondées sur les
principes de \textit{l'Information Processing}, de l'input compréhensible et de l'immersion. Ces
méthodes visent à construire la compétence procédurale avant d'exiger l'analyse déclarative.



\subsection{La lecture extensive (\textit{extensive reading})}

Pour contrer l'effet de surcharge cognitive, la recherche préconise la Lecture Extensive
(ER).\cite{II-3-1-ref23}

\begin{itemize}\item \textbf{Principe :} Lire de grandes quantités de textes faciles et compréhensibles, sans dictionnaire, pour le plaisir ou l'intérêt du contenu.\item \textbf{Mécanisme :} En lisant des textes où 95-98\% des mots sont connus (textes adaptés, « graded readers »), l'élève consolide la reconnaissance visuelle des mots (\textit{sight words}) et automatise le traitement syntaxique.\cite{II-3-1-ref21}\item \textbf{Application au Grec :} Cela implique de renoncer temporairement au « dogme des textes authentiques » pour utiliser des textes simplifiés ou construits pour l'apprentissage, permettant à l'élève de « lire » vraiment du grec dès la première année, au lieu de le déchiffrer.\cite{II-3-1-ref56}\end{itemize}

\subsection{L'immersion et la méthode polis}

La Méthode Polis, développée par Christophe Rico à Jérusalem, incarne l'application la plus aboutie
des principes neurocognitifs à l'enseignement des langues anciennes.\cite{II-3-1-ref58} Elle traite
le grec ancien comme une langue vivante (koinè), favorisant l'immersion totale.



\subsection{A. total physical response (tpr)}

Inspirée par James Asher, cette technique utilise l'impératif et le mouvement physique pour
enseigner le vocabulaire et les structures de base.\cite{II-3-1-ref61}

\begin{itemize}\item \textbf{Exemple :} L'enseignant dit « \textit{Boson\!} » (Lève-toi\!) et les élèves se lèvent.\item \textbf{Impact Cognitif :} Le TPR connecte directement le mot au geste et au concept, sans passer par la traduction en L1. Il ancre le langage dans la mémoire motrice et procédurale, contournant le filtre analytique de la mémoire déclarative.\cite{II-3-1-ref64}\end{itemize}

\subsection{B. living sequential expression (lse)}

Inspirée par François Gouin, cette technique propre à Polis consiste à décrire des séquences
d'actions logiques (« Je me lève, je vais à la porte, j'ouvre la porte ») en les mimant et en les
verbalisant.\cite{II-3-1-ref61}

\begin{itemize}\item \textbf{Impact Cognitif :} Elle respecte la séquentialité naturelle de l'action et permet d'assimiler les variations morphologiques (1ère personne, 3ème personne) par l'usage répété en contexte (\textit{usage-based learning}), favorisant l'induction naturelle des règles plutôt que leur déduction abstraite.\cite{II-3-1-ref61}\end{itemize}Ces méthodes « actives » ne sont pas de simples animations ludiques ; elles sont des dispositifs
rigoureux de \textbf{réduction de la charge cognitive extrinsèque} (pas de dictionnaire, pas
d'analyse méta-linguistique prématurée) pour maximiser la charge germane (acquisition du sens et des
structures).\cite{II-3-1-ref66}



\subsection{Repenser la place de la version}

Cette critique radicale de la version ne vise pas sa disparition, mais sa relocalisation
curriculaire.

Comme le suggère le Rapport Charvet, la version doit redevenir ce qu'elle aurait toujours dû être :
un exercice d'excellence pour des étudiants avancés, qui savent déjà lire le grec, et qui
travaillent la finesse de la transposition stylistique vers le français.\cite{II-3-1-ref16}

La séquence pédagogique doit être inversée :

1. \textbf{Phase d'Acquisition (Input/Immersion/Lecture Extensive) :} Construire le lecteur grec via
la mémoire procédurale.

2. \textbf{Phase d'Analyse (Grammaire explicite) :} Consolider les acquis par la réflexion
métalinguistique (mémoire déclarative).

3. \textbf{Phase de Traduction (Version) :} Exercer l'art de la correspondance entre deux systèmes
linguistiques maîtrisés.

Vouloir commencer par la version, c'est demander à l'élève de courir (traduire) avant de savoir
marcher (lire).

L'analyse approfondie des données historiques, cognitives et neuroscientifiques permet de confirmer
que l'assimilation de la « Version » à la « Lecture » constitue un mythe fondateur et délétère de
l'enseignement du grec ancien en France.

Ce mythe repose sur une méconnaissance profonde de la nature de l'acte de lire.

\begin{itemize}\item \textbf{Lire} est un processus de traitement de l'information (\textit{Information Processing}) qui exige rapidité, automaticité et fluidité, et qui s'appuie sur la mémoire procédurale.\item \textbf{Traduire (faire une version)} est un processus de résolution de problèmes (\textit{Problem Solving}) qui exige arrêt, analyse consciente et retour sur soi, et qui sollicite la mémoire déclarative au prix d'une charge cognitive massive.\end{itemize}En imposant la version comme unique voie d'accès au texte, l'institution scolaire a historiquement
privilégié la « gymnastique mentale » au détriment de la compétence linguistique réelle. Elle a
formé des générations d'élèves capables d'analyser grammaticalement une phrase de Xénophon sans
pouvoir la comprendre intuitivement.

Les sciences cognitives (LaBerge \& Samuels, Sweller) et les neurosciences (Ullman) nous enseignent
que cette approche est neurobiologiquement contre-productive pour l'apprentissage d'une langue. Pour
sauver l'enseignement des Humanités, il est impératif de dissocier l'apprentissage de la lecture de
l'exercice de traduction. L'adoption de méthodes basées sur l'input compréhensible, l'immersion
(Polis, TPR, LSE) et la lecture extensive n'est pas une concession à la facilité, mais un alignement
nécessaire sur le fonctionnement réel du cerveau humain. C'est à ce prix que le texte grec pourra
cesser d'être un code mort à déchiffrer pour redevenir une parole vivante à entendre.

